\RequirePackage{fix-cm}
\documentclass[11pt]{article}
\usepackage[margin=1in]{geometry}
\usepackage{fancyhdr}
\usepackage{enumitem}
\usepackage{parskip}
\usepackage{xcolor}

% Personal notes: \comment{Write your thoughts here.}
% Notes appear in grey in the compiled document; blank lines are supported.
\definecolor{commentgray}{gray}{0.5}
\newcommand{\comment}[1]{{\color{commentgray}#1}}
\usepackage{graphicx}
\usepackage{needspace}

\pagestyle{fancy}
\fancyhf{}
\fancyhead[C]{\textbf{SEG2105}}
\fancyfoot[C]{\thepage}
\setlength{\headheight}{14pt}
\setlength{\emergencystretch}{2em}

\newcounter{chapterNumber}
\setcounter{chapterNumber}{2}
\newcounter{headerOneCount}
\newcounter{headerTwoCount}[headerOneCount]
\setcounter{headerOneCount}{-1}
\renewcommand{\theheaderOneCount}{\arabic{headerOneCount}}
\renewcommand{\theheaderTwoCount}{\arabic{headerTwoCount}}

\newcommand{\Title}[1]{%
  \begin{center}
    \fontsize{16}{28}\selectfont\textbf{#1}
  \end{center}
}

\newcommand{\HeaderOne}[1]{%
  \Needspace{5\baselineskip}
  \refstepcounter{headerOneCount}%
  \setcounter{headerTwoCount}{-1}%
  {\noindent\fontsize{14}{22}\selectfont\textbf{\arabic{chapterNumber}.\theheaderOneCount\ #1}\par}
}

\newcommand{\HeaderTwo}[1]{%
  \Needspace{5\baselineskip}
  \refstepcounter{headerTwoCount}%
  {\noindent\fontsize{14}{20}\selectfont\textbf{\arabic{chapterNumber}.\theheaderOneCount.\theheaderTwoCount\ #1}\par}
}

\newcommand{\BodyText}[1]{%
  \fontsize{11}{15}\selectfont #1\par

}

\newenvironment{NumberedList}{%
  \begin{enumerate}[leftmargin=2em, itemsep=0.35em, topsep=0.35em]
}{%
  \end{enumerate}
}

\newenvironment{DottedList}{%
  \begin{itemize}[leftmargin=2em, itemsep=0.35em, topsep=0.35em, label=\textbullet]
}{%
  \end{itemize}
}

\begin{document}
\Title{Chapter 2: Review of Object Orientation}
\section*{Definitions}
\subsection*{Object-oriented foundations}
\begin{DottedList}
  \item \textbf{Abstraction:} Emphasizing relevant characteristics while leaving out unnecessary detail.
  \item \textbf{Procedural abstraction:} Representing a computation as a named procedure with a defined purpose and interface.
  \item \textbf{Data abstraction:} Representing information about an entity through a defined data model and operations while hiding unnecessary representation details.
  \item \textbf{Object orientation:} Organizing software as collaborating objects that combine data with the behaviour operating on it.
  \item \textbf{Modularity:} Dividing software into manageable units with clear responsibilities and interfaces.
  \item \textbf{Encapsulation:} Grouping state and behaviour within a unit such as a class and controlling access to it.
  \item \textbf{Information hiding:} Concealing implementation decisions behind an interface so other code need not depend on those details.
  \item \textbf{Object / instance:} An individual runtime entity with identity, state, and behaviour, created according to a class.
  \item \textbf{Class:} A description of the common structure and behaviour of a set of objects.
  \item \textbf{Identity:} What distinguishes an object from other objects, even when their data values are equal.
  \item \textbf{State:} The values and relationships held by an object at a particular time.
  \item \textbf{Behaviour:} The actions an object can perform and how it responds to operations.
  \item \textbf{Instance variable:} A field holding information separately for each instance of a class.
  \item \textbf{Attribute:} A property describing an object, commonly represented by a simple data value.
  \item \textbf{Association:} A relationship between objects.
  \item \textbf{Reference variable:} A variable whose value refers to an object; the variable is not the object itself.
  \item \textbf{Declared type:} The type specified for a variable, which determines the operations available through that variable at compile time.
  \item \textbf{Runtime class:} The actual class of an object during execution.
  \item \textbf{Class / static variable:} A field belonging to a class rather than separately to each instance.
  \item \textbf{Constant:} A named value intended to remain unchanged. A final reference cannot be reassigned, but its referenced object may still be mutable.
  \item \textbf{Mutable state:} Information that can change during execution.
\end{DottedList}
\subsection*{Methods, inheritance, and polymorphism}
\begin{DottedList}
  \item \textbf{Operation:} A description of behaviour independently of a particular implementation.
  \item \textbf{Method:} A class's implementation of an operation.
  \item \textbf{Polymorphism:} The ability to invoke the same operation on different kinds of objects and obtain the implementation appropriate to each object.
  \item \textbf{Inheritance:} A mechanism through which a subclass receives applicable features from a superclass and can add or specialize behaviour.
  \item \textbf{Superclass:} A more general class from which another class inherits.
  \item \textbf{Subclass:} A more specialized class that inherits from another class.
  \item \textbf{Inheritance hierarchy:} The organization of classes through superclass-subclass relationships.
  \item \textbf{Is-a rule:} The requirement that a subclass genuinely represents a kind of its superclass and supports the superclass contract.
  \item \textbf{Part-of relationship:} A relationship in which one object forms part of another, represented through associations rather than inheritance alone.
  \item \textbf{Contract:} The promised behaviour and conditions that users of a class or operation may rely on.
  \item \textbf{Abstract method:} A method declared without an implementation that concrete subclasses must implement or inherit an implementation for.
  \item \textbf{Abstract class:} A class that cannot be instantiated directly and may leave some operations unimplemented.
  \item \textbf{Concrete class:} A class that can be instantiated and provides implementations for all required abstract methods.
  \item \textbf{Overriding:} Supplying a subclass implementation of an inherited instance method.
  \item \textbf{Overloading:} Defining methods with the same name but different parameter lists.
  \item \textbf{Dynamic binding:} Selecting an overridden instance method at runtime according to the actual class of the receiving object.
  \item \textbf{Cast:} An explicit conversion of the type through which a value is viewed. Casting a reference does not change the object's class.
  \item \textbf{Downcast:} Casting a reference to a more specific type; an incompatible downcast fails at runtime.
\end{DottedList}
\subsection*{Java language and collections}
\begin{DottedList}
  \item \textbf{Bytecode:} Instructions produced by the Java compiler for execution by a Java virtual machine.
  \item \textbf{Java virtual machine (JVM):} The execution environment that runs Java bytecode.
  \item \textbf{API:} The accessible types and operations through which code uses a software component.
  \item \textbf{Javadoc:} A tool and documentation-comment format for generating API documentation from Java source code.
  \item \textbf{\texttt{char}:} A Java primitive value representing one UTF-16 code unit.
  \item \textbf{Unicode:} A standard assigning code points to characters used in writing systems and symbols.
  \item \textbf{UTF-16 code unit:} A 16-bit unit used to encode Unicode text; some characters require two such units.
  \item \textbf{Wrapper class:} An object type representing a primitive value, such as \texttt{Character} for \texttt{char}.
  \item \textbf{\texttt{String}:} Java's immutable representation of a sequence of text characters.
  \item \textbf{Concatenation:} Joining text sequences together, for example using \texttt{+} with strings.
  \item \textbf{Array:} An indexed container with a fixed length and a declared element type.
  \item \textbf{Collection:} An object that groups elements; Java's \texttt{Collection} interface defines common collection operations.
  \item \textbf{\texttt{ArrayList}:} A resizable, indexed list implementation in Java.
  \item \textbf{Iterator:} An object used to traverse elements, commonly through \texttt{hasNext()} and \texttt{next()}.
  \item \textbf{Generics:} Type parameters that express the types a class or method works with, improving compile-time type checking.
  \item \textbf{Exception:} An object signalling a failure or abnormal condition that interrupts normal execution flow.
  \item \textbf{\texttt{try}/\texttt{catch}:} Java constructs that run code and handle exceptions of specified types.
  \item \textbf{Interface:} A contract that implementing classes promise to support; variables can use an interface type to refer to different implementations.
  \item \textbf{Default interface method:} An instance method with an implementation supplied by an interface.
  \item \textbf{Static method:} A method associated with a class or interface rather than invoked through instance-method dynamic binding.
  \item \textbf{\texttt{Comparable}:} An interface for defining an object's natural ordering through comparison with another object.
\end{DottedList}
\subsection*{Organization, access, and execution}
\begin{DottedList}
  \item \textbf{Package:} A grouping of related Java types that provides a namespace and an access boundary.
  \item \textbf{Namespace:} A naming context that distinguishes otherwise identical names.
  \item \textbf{Import:} A declaration allowing a type or static member to be referenced by a shorter name in source code.
  \item \textbf{Visibility / access control:} Rules determining which code may use a type or member.
  \item \textbf{\texttt{public}:} An access modifier permitting broad access, subject to the accessibility of the enclosing type and module.
  \item \textbf{\texttt{private}:} An access modifier limiting access to the declaring class and its enclosing nest of related nested classes.
  \item \textbf{Package access:} Access available within the same package when no access modifier is specified.
  \item \textbf{\texttt{protected}:} Access within the package and from subclasses, with additional restrictions for subclass access across packages.
  \item \textbf{Thread:} A sequence of execution that can progress concurrently with other threads.
  \item \textbf{Concurrency:} Multiple activities making progress during overlapping periods, whether or not they execute simultaneously.
  \item \textbf{\texttt{Runnable}:} An interface defining a \texttt{run()} operation for a task that can be executed by a thread.
  \item \textbf{\texttt{Thread.start()}:} An operation that starts a separate thread which executes its task's \texttt{run()} method.
  \item \textbf{\texttt{run()}:} The method containing a thread task's work; calling it directly executes on the calling thread.
  \item \textbf{User interface (UI):} The part of a system through which users provide input and receive output.
  \item \textbf{Domain logic:} Code implementing the rules and operations of the problem domain.
  \item \textbf{Deprecated feature:} A feature marked as discouraged for continued use, often because a preferred replacement exists.
\end{DottedList}
\newpage

\HeaderOne{What is Object Orientation?}
\begin{DottedList}
  \item Procedural abstraction organizes computation into procedures. Data abstraction groups information describing an entity.
  \item Object orientation places behaviour with the data it operates on. A running program consists of collaborating objects.
  \item Abstraction emphasizes relevant characteristics. Modularity divides code into manageable units. Encapsulation groups state and behaviour and supports information hiding.
\end{DottedList}
\HeaderOne{Classes and Objects}
\begin{DottedList}
  \item An object is an individual runtime entity with identity, state, and behaviour. A class describes its instances' common structure and behaviour.
  \item Film can be a class and a particular film an instance. Genre can be a class and Science Fiction one instance.
  \item Choose singular, descriptive class names such as BankAccount, with the right level of generality and one clear meaning.
\end{DottedList}
\HeaderOne{Instance Variables}
\begin{DottedList}
  \item Instance variables hold information for each object. Attributes hold simple values; associations represent relationships with other objects.
  \item A reference variable refers to an object; it is not the object itself. Multiple variables can refer to the same object, and a variable may later refer to another object.
  \item The variable's declared type determines which operations are available through it.
\end{DottedList}
\HeaderTwo{Class variables}
\begin{DottedList}
  \item A static variable belongs to the class and is shared among its instances.
  \item Constants and lookup data are common uses. Mutable shared state needs care because a change affects all its users.
\end{DottedList}
\HeaderOne{Methods, Operations, and Polymorphism}
\begin{DottedList}
  \item An operation describes behaviour independently of its implementation. A method supplies a class's implementation.
  \item Polymorphism allows the same operation to behave differently for different kinds of objects; circles and rectangles calculate area differently.
  \item Polymorphic methods reduce repeated type tests and switch statements.
\end{DottedList}
\HeaderOne{Inheritance Hierarchies}
\begin{DottedList}
  \item A superclass defines common features. Subclasses inherit applicable features and may add or specialize behaviour.
  \item Apply the is-a rule: a checking account is an account, but a province is not a country.
  \item Every inherited feature must make sense in the subclass. A part-of relationship is an association rather than inheritance.
\end{DottedList}
\HeaderOne{Inheritance, Polymorphism, and Variables}
\HeaderTwo{Abstract classes and overriding}
\begin{DottedList}
  \item Declare an operation at the highest sensible level. An abstract method specifies behaviour without implementing it there.
  \item An abstract class cannot be instantiated. A concrete class must implement or inherit implementations of all required abstract methods.
  \item Overriding supplies a subclass implementation of an inherited instance method. It can extend or optimize behaviour while preserving the contract.
  \item Overloading uses different parameter lists and is distinct from overriding.
\end{DottedList}
\HeaderTwo{Dynamic binding}
\begin{DottedList}
  \item The declared type controls whether a call is legal; the runtime object's class determines which overridden instance method executes.
  \item Look first for an implementation in the runtime class, then up its superclass chain.
  \item Example: Shape s = new Circle(); s.area(); executes the Circle implementation.
  \item Casting changes the type through which a reference is viewed, not the object's class. An invalid downcast can fail at runtime.
\end{DottedList}
\HeaderOne{Java Fundamentals}
\HeaderTwo{Language and documentation}
\begin{DottedList}
  \item The lecture traces object orientation through Simula, Smalltalk, and C++, then Java's development from Oak and its 1995 launch.
  \item Java compiles to bytecode executed by a virtual machine. Read class and method documentation to understand library behaviour.
  \item Javadoc generates API documentation from source declarations and comments.
\end{DottedList}
\HeaderTwo{Characters, strings, and collections}
\begin{DottedList}
  \item char holds one UTF-16 code unit; Character is its wrapper class. Some Unicode characters require two code units.
  \item String represents text; + concatenates strings.
  \item Arrays have fixed length. ArrayList supports a resizable collection. Iterators traverse elements using hasNext() and next().
  \item Generics express element types: ArrayList\textless{}String\textgreater{} names = new ArrayList\textless{}\textgreater{}(); avoids many explicit casts during retrieval.
\end{DottedList}
\HeaderTwo{Exceptions and interfaces}
\begin{DottedList}
  \item Exceptions signal failures. A try/catch block handles specified types of exceptions; the handling should reflect the actual problem.
  \item An interface specifies behaviour implementing classes promise to provide. A class can implement multiple interfaces, and a variable can have an interface type.
  \item Runnable, Collection, Iterator, and Comparable are library examples.
  \item The slides' prohibition on implementation code in interfaces describes older Java. Modern Java permits default and static interface methods.
\end{DottedList}
\HeaderTwo{Packages and access}
\begin{DottedList}
  \item Packages group related classes and provide namespaces. Imports allow short type names in source code.
  \item public permits broad access; private hides implementation within the class and its enclosing nest; no modifier gives package access.
  \item protected permits package access and subclass access, subject to Java's additional rules across packages.
  \item Use the narrowest practical visibility.
\end{DottedList}
\HeaderTwo{Threads}
\begin{DottedList}
  \item A thread is an execution sequence that can progress concurrently with other threads.
  \item Implement Runnable.run(), create a Thread using that Runnable, and call start(). Alternatively, subclass Thread.
  \item Calling run() directly uses the calling thread; start() begins a separate thread of execution.
  \item Coordinate access when threads share mutable state.
\end{DottedList}
\HeaderOne{Programming Style}
\begin{DottedList}
  \item Prefer understandable code and descriptive names. Shorter or cleverer code is not automatically better.
  \item Explain non-obvious decisions in comments and organize class elements consistently.
  \item Extract shared behaviour into methods instead of copying code, so fixes do not need to be repeated.
  \item Follow the is-a rule, hide implementation details, and separate UI code from domain logic to improve reuse.
\end{DottedList}
\HeaderOne{Difficulties and Risks}
\begin{DottedList}
  \item Languages and libraries evolve. Deprecated features may need replacement during maintenance.
  \item Measure efficiency against the application's requirements. Virtual-machine execution and dynamic binding alone do not establish that a program is too slow.
\end{DottedList}
\end{document}
